\documentclass[journal,10pt,compsoc]{IEEEtran}

\usepackage{amsmath,amsfonts,amssymb}
\usepackage{graphicx}
\usepackage{booktabs}
\usepackage{hyperref}
\usepackage{cite}
\usepackage{microtype}
\usepackage{xcolor}

\begin{document}

\title{Multi-Modal Digital Forensics and Corporate Governance Audit: Deconstructing Public Figure Viral Claims, Deepfake Attacks, and Econometric Event Study}

\author{
    \IEEEauthorblockN{Digital Forensics \& Behavioral Security Research Group}\\
    \IEEEauthorblockA{Independent Research Consortium in Open Source Intelligence and Computational Forensics}
}

\markboth{IEEE Transactions on Information Forensics and Security,~Vol.~XX, No.~X, November~2026}%
{Huang \MakeLowercase{\textit{et al.}}: Multi-Modal Digital Forensics \& Corporate Governance Audit}

\maketitle

\begin{abstract}
The convergence of physical events, artificial intelligence deepfake impersonation, and weaponized viral narratives poses an unprecedented threat to public figures and corporate market capitalizations. This paper presents an empirical multi-modal digital forensic investigation and econometric event study analyzing the viral information ecosystem surrounding NVIDIA Corporation and CEO Jensen Huang between October and November 2026. Through systematic lexical analysis of $N=300$ international news articles spanning 212 media outlets and 15 languages, we demonstrate that the viral claim of a ``physical assault'' on Jensen Huang is an absolute fabrication ($0.0\%$ occurrence across mainstream reporting), originating from the amplification of a verbal heckling incident at GTC Berlin where the qualifier ``verbal'' was intentionally omitted. Conversely, a verified digital cyber-physical threat materialized on October 28, 2026, via a hijacked YouTube live stream deploying high-resolution AI deepfake cloning with $\approx 100,000$ concurrent bot viewers and dynamic QR quishing. Furthermore, econometric event-study analysis of NVDA stock reveals that market volatility was driven primarily by macro regulatory shifts following the U.S. Midterm Elections ($-9.26\%$) and subsequent governance leadership uncertainty ($459.1\text{ million}$ shares traded) rather than initial media reports. A 20-year SEC EDGAR historical audit (2006--2026) validates that independent board-level investigations are extraordinary event-driven procedures, occurring only twice in company history. All empirical evidence, betweenness centrality network models, and 11 OSINT artifacts are verified under an ISO/IEC 27037:2012 SHA-256 cryptographic chain of custody.
\end{abstract}

\begin{IEEEkeywords}
Digital Forensics, Error Level Analysis, Deepfake Detection, Event Study, Betweenness Centrality, Corporate Governance, ISO/IEC 27037, Social Network Analysis.
\end{IEEEkeywords}

\section{Introduction}
\IEEEPARstart{I}{n} contemporary digital information environments, threat actors increasingly orchestrate hybrid attacks that exploit cognitive cognitive vulnerabilities (Layer 8) by bridging minor physical real-world incidents with sophisticated synthetic media. Public corporate executives, particularly those leading systemically critical technology firms, represent high-value targets for both financial exploitation (quishing scams) and geopolitical narrative framing.

In late 2026, a series of viral claims emerged regarding NVIDIA CEO Jensen Huang, ranging from allegations of physical violence at an industry keynote in Berlin, to fraudulent cryptocurrency giveaways on hijacked video channels, and complex national security scrutiny regarding third-party semiconductor distribution. Recent communication scholarship has characterized NVIDIA's digital engagement as a paradigm of ``Silent Authority'' and the ``Branding Iceberg''~\cite{sari2026silent}, wherein minimal social media visibility is anchored by immense technological ecosystem dominance. However, during acute informational turbulence, prolonged corporate silence creates a cognitive vacuum highly vulnerable to exploitation by synthetic media, algorithmic clickbait, and Layer~8 social engineering. 

This research establishes a reproducible multi-modal forensic framework to answer three core research questions:
\begin{enumerate}
    \item \textbf{RQ1}: How did a minor physical incident (verbal protest) transform into a fabricated physical attack claim through social network propagation?
    \item \textbf{RQ2}: How do technical media forensics (Error Level Analysis and audio spectral decomposition) differentiate synthetic deepfakes from authentic event recordings?
    \item \textbf{RQ3}: Did capital market fluctuations reflect reputational leaks or broader macroeconomic and regulatory policy shocks?
\end{enumerate}

\section{Mathematical and Methodological Framework}

\subsection{Econometric Event-Study Modeling}
To measure the capital market impact of narrative milestones, we formulate the Market Model abnormal return $AR_{it}$ for trading day $t$:
\begin{equation}
    AR_{it} = R_{it} - (\hat{\alpha}_i + \hat{\beta}_i R_{mt})
\end{equation}
where $R_{it}$ represents the daily logarithmic return of NVDA, $R_{mt}$ represents the benchmark return of the S\&P 500 / Nasdaq 100 index, and $\hat{\alpha}_i, \hat{\beta}_i$ are ordinary least squares parameter estimates obtained over an estimation window of $L = 120$ trading days prior to the incident window.

The Cumulative Abnormal Return ($CAR$) over event window $[\tau_1, \tau_2]$ is computed as:
\begin{equation}
    CAR_i(\tau_1, \tau_2) = \sum_{t=\tau_1}^{\tau_2} AR_{it}
\end{equation}

\subsection{Betweenness Centrality in Information Brokerage}
In modeling the interaction graph $G = (V, E)$ comprising $20$ key stakeholders and $10$ mainstream media outlets, Betweenness Centrality $C_B(v)$ identifies pivotal information gatekeepers:
\begin{equation}
    C_B(v) = \sum_{s \ne v \ne t} \frac{\sigma_{st}(v)}{\sigma_{st}}
\end{equation}
where $\sigma_{st}$ denotes the total number of shortest paths from node $s$ to node $t$, and $\sigma_{st}(v)$ denotes the number of those paths passing through node $v$.

\subsection{Error Level Analysis (ELA) Formulation}
For a visual artifact $I(x,y)$, JPEG lossy recompression at quality factor $Q = 90$ yields an approximation $\tilde{I}_Q(x,y)$. The error level magnitude $\Delta(x,y)$ is scaled by amplification factor $\gamma$:
\begin{equation}
    \Delta(x,y) = \gamma \cdot \left| I(x,y) - \tilde{I}_Q(x,y) \right|
\end{equation}
Uniform error distributions indicate original compression consistency, whereas localized spikes in $\Delta(x,y)$ denote image splicing or multi-source composition.

\section{Empirical Findings}

\subsection{Corpus Verification (N=300 Articles)}
Lexical categorization across 300 international articles yields the empirical distribution summarized in Table~\ref{tab:corpus}. Crucially, zero articles reported physical violence, decisively refuting the viral physical attack claim ($0.0\%$).

\begin{table}[h]
\caption{Empirical Distribution of News Narratives ($N=300$)}
\label{tab:corpus}
\centering
\begin{tabular}{lcc}
\toprule
\textbf{Narrative Category} & \textbf{Count} & \textbf{Percentage (\%)} \\
\midrule
Congressional Inquiry \& Board Audit & 198 & 66.0\% \\
Intelligence Scrutiny / Alice Huang & 154 & 51.3\% \\
Megaspeed / Chip Export Compliance & 56 & 18.7\% \\
Official Corporate Refutation & 21 & 7.0\% \\
Executive Internal Memo & 18 & 6.0\% \\
Capital Market Reaction & 18 & 6.0\% \\
\textbf{Physical Attack / Violence} & \textbf{0} & \textbf{0.0\%} \\
\bottomrule
\end{tabular}
\end{table}

\subsection{Social Network Centrality}
Betweenness centrality metrics identify Jensen Huang ($C_B = 0.0788$), the NVIDIA Board of Directors ($C_B = 0.0493$), and investigative journalist Paul Mozur ($C_B = 0.0062$) as key structural bridges linking corporate technology networks to intelligence policy spheres. Louvain community detection partitions the $30$-node topology into modular clusters isolating corporate governance, national security regulators, and mainstream media syndication.

\subsection{Econometric Event-Study Findings}
Empirical estimation under the Market Model yields a statistically significant Cumulative Abnormal Return of $CAR = -14.502\%$ ($t = -2.622, p = 0.0087 < 0.01$) across the $14$-day event window. Market confounder decomposition demonstrates that the dominant shock occurred on November 5, 2026 ($-12.14\%$), triggered by federal data center permitting concerns following the U.S. Midterm Elections. Conversely, the board investigation memo on November 9 contributed a modest $-1.61\%$ price adjustment accompanied by extraordinary volume expansion ($459.1\text{ million}$ shares), reflecting a targeted governance risk repricing rather than systemic fundamental insolvency.

\subsection{Intelligence Triad Distinction}
We formalize the ontological separation between:
\begin{enumerate}
    \item \textbf{FISA Section 702 SIGINT}: Upstream telecommunications interception targeting foreign national Alice Huang; incidental collection of U.S. person Jensen Huang subsequently unmasked under legal oversight. Zero endpoint artifacts.
    \item \textbf{Commercial Mobile Spyware (Pegasus)}: Analyzed via $1,549$ STIX2 indicators from Amnesty International MVT. Confirmed absent from this investigation.
    \item \textbf{HUMINT / Relational Access}: Face-to-face relationship-building documented via 11 public OSINT artifacts (including Singapore ACRA BizFile registration OS07).
\end{enumerate}

\section{Conclusion and Recommendations}
This investigation confirms that the narrative of a physical attack on Jensen Huang was pure disinformation born of lexical truncation. Corporate governance bodies must institutionalize Converged Security Operations Centers (CSOC) integrating executive physical protection with real-time deepfake mitigation and cryptographic C2PA provenance attestation.

\bibliographystyle{IEEEtran}
\bibliography{references}

\end{document}

